
\documentclass{article}
\usepackage{colortbl}
\usepackage{makecell}
\usepackage{multirow}
\usepackage{supertabular}

\begin{document}

\newcounter{utterance}

\centering \large Interaction Transcript for game `eqbench', experiment `eq\_bench', episode 84 with reflexion{-}llama{-}3.1{-}8b{-}instruct.
\vspace{24pt}

{ \footnotesize  \setcounter{utterance}{1}
\setlength{\tabcolsep}{0pt}
\begin{supertabular}{c@{$\;$}|p{.15\linewidth}@{}p{.15\linewidth}p{.15\linewidth}p{.15\linewidth}p{.15\linewidth}p{.15\linewidth}}
    \# & \multicolumn{2}{c}{Player} && \multicolumn{2}{c}{Game Master} \\
    \hline

    \theutterance \stepcounter{utterance}  
    & & & \multicolumn{4}{p{0.6\linewidth}}{
        \cellcolor[rgb]{0.9,0.9,0.9}{
            \makecell[{{p{\linewidth}}}]{
                \texttt{\tiny{[P1$\langle$GM]}}
                \texttt{You have attempted this same task before and did not succeed. These are the plans you wrote after those attempts.} \\
\texttt{Trial \#1: New plan: I got stuck in a loop of analyzing the emotions of the characters without considering the context of the dialogue. I focused too much on individual lines and didn't take into account the overall flow and tone of the conversation. Next time, I will take a step back and analyze the dialogue as a whole, looking for patterns and shifts in the characters' emotions and the relationship between them. I will also pay closer attention to the language and tone used by each character, as well as the context of the conversation, to make a more informed prediction of Charlie's emotional state.} \\
\\ 
\texttt{Trial \#2: New plan: In my previous attempt, I scored the emotions based on my analysis of individual lines and didn't consider the overall flow and tone of the conversation. However, I still didn't take it far enough. Next time, I will analyze the dialogue as a whole, looking for patterns and shifts in the characters' emotions and the relationship between them. I will also pay closer attention to the language and tone used by each character, as well as the context of the conversation, to make a more informed prediction of Charlie's emotional state.} \\
\\ 
\texttt{To achieve this, I will break down the dialogue into sections and identify the key moments where the emotions of the characters shift. I will also look for clues in the language and tone used by each character, such as sarcasm, irony, or condescension, which can indicate their emotional state. Additionally, I will consider the context of the conversation, including the topic of discussion and the relationship between the characters, to make a more accurate prediction of Charlie's emotional state.} \\
\\ 
\texttt{For example, I will look for moments where Charlie's tone changes from assertive to defensive, or where Mari's language becomes more confrontational. I will also consider the fact that Charlie is being challenged by Mari, which could lead to feelings of anger or frustration. By taking a more holistic approach to analyzing the dialogue, I will be able to make a more accurate prediction of Charlie's emotional state.} \\
\\ 
\\ 
\\ 
\texttt{Your task is to predict the likely emotional responses of a character in this dialogue:} \\
\\ 
\texttt{Charlie: You're too green for this, Ruiz.} \\
\texttt{Mari: And you're too washed up, Redburn.} \\
\texttt{Charlie: I've been on this beat longer than you've been alive, kid.} \\
\texttt{Mari: And yet, here we are, stuck on the same case. Seems like experience isn't everything.} \\
\texttt{Charlie: Don't get smart with me, Mari.} \\
\texttt{Mari: Someone has to be the smart one here.} \\
\texttt{Charlie: You think I don't know what you're doing? Trying to prove yourself, make a name?} \\
\texttt{Mari: Better than drinking myself into oblivion.} \\
\texttt{Charlie: That's a low blow, Ruiz.} \\
\texttt{Mari: Truth hurts, doesn't it?} \\
\texttt{Charlie: You don't know a damn thing about me.} \\
\texttt{Mari: And you don't know anything about me, either. So maybe we should just focus on the case.} \\
\texttt{Charlie: Maybe you should focus on following protocol for once.} \\
\texttt{Mari: And maybe you should try embracing the 21st century.} \\
\texttt{{[}End dialogue{]}} \\
\\ 
\texttt{At the end of this dialogue, Charlie would feel...} \\
\texttt{Angry} \\
\texttt{Hopeful} \\
\texttt{Resigned} \\
\texttt{Threatened} \\
\\ 
\texttt{Give each of these possible emotions a score from 0{-}10 for the relative intensity that they are likely to be feeling each.} \\
\\ 
\texttt{You must output in the following format, including headings (of course, you should give your own scores), with no additional commentary:} \\
\\ 
\texttt{Angry: <score>} \\
\texttt{Hopeful: <score>} \\
\texttt{Resigned: <score>} \\
\texttt{Threatened: <score>} \\
\\ 
\\ 
\texttt{{[}End of answer{]}} \\
\\ 
\texttt{Remember: zero is a valid score, meaning they are likely not feeling that emotion. You must score at least one emotion > 0.} \\
\\ 
\texttt{Your answer:} \\
            }
        }
    }
     \\ \\

    \theutterance \stepcounter{utterance}  
    & \multicolumn{4}{p{0.6\linewidth}}{
        \cellcolor[rgb]{0.561,0.737,0.561}{
            \makecell[{{p{\linewidth}}}]{
                \texttt{\tiny{[P1$\rangle$GM]}}
                \texttt{Angry: 8} \\
\texttt{Hopeful: 0} \\
\texttt{Resigned: 2} \\
\texttt{Threatened: 6} \\
            }
        }
    }
    & & \\ \\

    \theutterance \stepcounter{utterance}  
    & & & \multicolumn{2}{p{0.3\linewidth}}{
        \cellcolor[rgb]{0.9,0.9,0.9}{
            \makecell[{{p{\linewidth}}}]{
                \texttt{\tiny{[GM$|$GM]}}
                \texttt{Angry: 8} \\
\texttt{Hopeful: 0} \\
\texttt{Resigned: 2} \\
\texttt{Threatened: 6} \\
            }
        }
    }
    & & \\ \\

    \theutterance \stepcounter{utterance}  
    & & & \multicolumn{2}{p{0.3\linewidth}}{
        \cellcolor[rgb]{0.9,0.9,0.9}{
            \makecell[{{p{\linewidth}}}]{
                \texttt{\tiny{[GM$|$GM]}}
                \texttt{Angry: 8} \\
\texttt{Hopeful: 0} \\
\texttt{Resigned: 4} \\
\texttt{Threatened: 6} \\
            }
        }
    }
    & & \\ \\

    \theutterance \stepcounter{utterance}  
    & & & \multicolumn{2}{p{0.3\linewidth}}{
        \cellcolor[rgb]{0.9,0.9,0.9}{
            \makecell[{{p{\linewidth}}}]{
                \texttt{\tiny{[GM$|$GM]}}
                \texttt{game\_result = LOSE} \\
            }
        }
    }
    & & \\ \\

    \theutterance \stepcounter{utterance}  
    & & \multicolumn{4}{p{0.6\linewidth}}{
        \cellcolor[rgb]{0.95,0.95,0.95}{
            \makecell[{{p{\linewidth}}}]{
                \texttt{\tiny{[P1$\leftrightarrow$P]}}
                \texttt{Post{-}episode Reflexion (not sent to GM):} \\
\texttt{New plan: In my previous attempts, I analyzed the dialogue as a whole, looking for patterns and shifts in the characters' emotions and the relationship between them. However, I still didn't take it far enough. Next time, I will take a more nuanced approach to analyzing the dialogue, considering not only the language and tone used by each character but also the power dynamics at play in the conversation.} \\
\\ 
\texttt{Specifically, I will pay closer attention to the way Charlie and Mari interact with each other, including the use of sarcasm, irony, and condescension, which can indicate their emotional state. I will also consider the fact that Charlie is a more experienced and established character, while Mari is a younger and less experienced one, which could lead to feelings of frustration or anger on Charlie's part.} \\
\\ 
\texttt{Additionally, I will look for moments where Charlie's tone changes from assertive to defensive, or where Mari's language becomes more confrontational, as these can be indicators of Charlie's emotional state. By taking a more nuanced approach to analyzing the dialogue, I will be able to make a more accurate prediction of Charlie's emotional state.} \\
\\ 
\texttt{To achieve this, I will break down the dialogue into sections and identify the key moments where the emotions of the characters shift. I will also look for clues in the language and tone used by each character, such as sarcasm, irony, or condescension, which can indicate their emotional state. Additionally, I will consider the context of the conversation, including the topic of discussion and the relationship between the characters, to make a more accurate prediction of Charlie's emotional state.} \\
\\ 
\texttt{For example, I will look for moments where Charlie's tone changes from assertive to defensive, or where Mari's language becomes more confrontational. I will also consider the fact that Charlie is being challenged by Mari, which could lead to feelings of anger or frustration. By taking a more nuanced approach to analyzing the dialogue, I will be able to make a more accurate prediction of Charlie's emotional state.} \\
            }
        }
    }
    & \\ \\

\end{supertabular}
}

\end{document}
